\scanpage{57}
\deleted{\unsure{Faisant|Passant}} Prenant pour $X,Y$ des VA, on trouve :

\uline{Corollaire} 1.7. \deleted{\unsure{Tout|Soit}} Soient $X,Y$ deux V.A. de dim. $n,m$. Pour qu'un homomorphisme $\psi:H^1(X)\longrightarrow H^1(Y)$ soit définissable par un cycle algébrique $Z$ de \deleted{\unsure{di|dim}} codim $n$ sur $X\times Y$, il faut et il suffit qu'il soit définissable par un homomorphisme de V.A. $u:Y\unsure{-|\longrightarrow}X$ comme $\psi=u^{(1)}$. De plus, $Z$ étant fixé, \deleted{\unsure{ce|u}} $u$ est uniquement déterminé, et indépendant du nb premier $l$ choisi pour les coefficients dans la cohomologie.

\uline{Remarque} 1.8. Généralisation \inserted{de 1.1.} aux schémas projectifs lisses, à fibres \unsure{géom|géom.} connexes, sur une base quelconque $S\deleted{\unsure{Z}}$, en partant d'un $n$-cocycle (au sens du \deleted{\unsure{le}} gradué associé au $K$) $Z\in{\mathrm{CK}}^n(X\times_S Y)$ : on obtient, si $A_Y$ et $A_X$ existent, (plus gén, si $A_Y$ existe
% archive sheet number: 56
